\chapter{研究实施与质量控制}
\chapterauthors{Deborah G. Grady 和 Alison J. Huang}
\sourcepage{561}
本书大部分讨论设计，本章着重实施与质量控制。所有临床研究都面临实施挑战，直接招募参与者收集资料的研究最为复杂。即使坐在书桌前深思熟虑的最佳计划，实际也可能不同：找不到熟练人员，场地不理想，参与者不愿入组，随访比预期困难，干预耐受性差，测量难以取得。无知、粗心、缺乏培训和标准化，以及方案定稿与实施中的其他错误，都会损害良好设计的结论。

成功实施始于汇集场地、设备、人员和财务管理等资源。随后通过预试招募机制、测量程序及适用的干预计划，使方案定稿，避免开始收集资料后再修改。\term{操作手册}应详列具体程序，研究开始后可能仍需更新；之后系统开展资料收集和管理的\term{质量控制}。

本章多数策略针对小团队单中心研究，但许多也适用于多位研究者领导、跨多个场所的大型团队项目。

\sourcepage{562}
\section{准备工作}
\subsection{研究启动}
开始时，\term{主要研究者}（principal investigator，PI）与人员需接受方案培训，确定预算、签订合同、设置岗位、招聘培训、配置场地、取得伦理审查批准、购买材料、编写操作手册、开发测试资料表和数据库，并规划招募策略和材料。首名参与者入组前这一阶段称为\term{研究启动}（study start-up），工作密集，常需数月。留足时间、做好规划，对高质量研究很重要。

\subsection{研究团队}
团队可小至研究者加兼职助理，也可有数百名全职人员，但活动相似，即使一人兼多职。PI 必须保证表18.1各项职能得到履行。财务、人力资源等人员，通常由机构聘用，为某单位全部研究者和工作人员提供支持。

\begingroup\small
\begin{longtable}{L{.20\textwidth}L{.36\textwidth}L{.35\textwidth}}
\caption*{\textbf{表18.1\quad 研究团队成员的职能角色}}\\\toprule
角色 & 职能 & 说明\\\midrule\endfirsthead\toprule 角色 & 职能 & 说明\\\midrule\endhead
主要研究者（PI） & 对设计、经费、人员、实施、质量、参与者安全和结果报告负最终责任。 & \\
项目主管／临床协调员 & 日常管理全部研究活动。 & 应有经验、负责、细致，具备良好人际与组织能力。\\
\sourcepage{563}招募协调员 & 确保达到合格参与者入组目标。 & 应熟悉多种招募技术。\\
临床研究协调员／临床人员 & 执行访视程序和测量。 & 体检或其他专门操作可能需特别执照或认证。\\
临床研究监察员 & 监测场所活动，尤其多中心研究。 & 现场访视、审查文件，确保进展适当、遵守监管要求并妥善记录。\\
监管事务经理／质控协调员 & 向伦理审查委员会和政府机构处理、提交文件，保证合规与时限；确保人员遵守标准操作规程，监督质控。 & 办理伦理和监管申请、维持批准；观察操作，确保按规程执行；可能负责配合 FDA 等外部审计。\\
数据管理员 & 设计、测试和实施录入、编辑与管理系统。 & 持续维护清理、报告、质疑和修改系统。\\
程序员／分析员 & 编制招募、依从性和数据质量报告，开展分析。 & 在 PI 与统计师监督下工作。\\
统计师 & 参与设计，估计样本量与效能，设计分析计划及数据安全监察指南，解释发现。 & 参与总体设计、实施、期中监测、分析和结果呈现。\\
\sourcepage{564}行政助理 & 提供行政支持、安排会议等。 & \\
财务经理 & 编预算、管理支出。 & 提供预测，协助管理预算。\\
人力资源经理 & 协助职位说明、招聘和考评。 & 帮助处理人事问题。\\\bottomrule
\end{longtable}\endgroup

确定人数和分工后，与部门行政人员拟定、发布职位说明，再评估申请、面试。招聘并不容易：某些职能无正式培训，不同项目要求也不同。项目主管负责日常管理，可有护理、药学、公共卫生、实验室或药物研究背景，职责差异很大。另一重要岗位是\term{临床研究协调员}（clinical research coordinator，CRC），执行日常程序、进行测量，通常与参与者接触最多，是研究的“面孔”。为做好招募、留存，CRC 必须建立融洽关系、回应需求。许多机构提供包括 CRC 在内的培训认证。

除正式招聘渠道，在线职位、社交媒体或 LinkedIn 也有用。较有把握的是选择能力已知者，如同事项目结束的人员。
\sourcepage{565}
有些机构备有经验丰富的协调员等，可兼职聘用。

\subsection{领导与团队建设}
\textbf{研究质量始于 PI 的诚信与领导。}PI 应确保所有人员受训并具备履职资格，明确传达：保护参与者和隐私、数据完整准确、公平呈现发现最为重要。虽不能监看每次测量，却可让团队感到负责人了解全部活动、重视保护与质量，多数人员会相应回应。好的领导者适当授权，同时建立监督体系。

规划伊始就应定期召集团队会议，预先发议程，让各领域负责人报告。会议既发现、解决问题，也让所有人参与项目发展。可加入科学讨论、初步发现介绍，增进同事情谊、士气和目标兴趣，并提供在岗教育培训。

还可能需定期一对一会谈，让人员提出不便在集体场合说的关切和改进建议。工作表现反馈应个别进行；虽费时费力，但研究者既是领导者，也是雇主，不能忽视。
\sourcepage{566}
许多机构提供数据库管理、专业招募、专门测量核心实验室、熟悉 FDA 等监管问题的分析人员，以及表格文件库。大型机构中获取渠道未必显而易见，应先熟悉当地资源，避免重复建设。

\subsection{场地与设备}
有些研究在专用或与临床共用的场地访视；有些依赖网络互动系统、邮寄干预、电话视频访视、居家测量、在线或移动录入。即使主要远程开展，人员仍需专用空间处理后台工作，存放、组织和加工材料。

若需实体场地，应便利、舒适、价格可承受且足够宽敞。未能尽早落实，可能导致招募困难、访视依从差、数据不全和人员不满。场地应容易找到，有停车位或临近公交，容纳人员、设备、研究药物及文件。体检需隐私和洗手条件；需去实验室或放射科时，也应方便。

许多研究需电脑、实验室仪器及体检设备。启动时可能需购买，并安排维护、校准和质量检查。
\sourcepage{567}
昂贵设备应留出多家报价和机构特别审批的时间。若有集中临床研究设施，使用它们可能更方便，但通常需付费并经过申请审查。

\subsection{研究预算}
充足资金至关重要。PI 早在申请资助时已编预算（第20章）。多数机构有资助前管理人员协助；应熟悉他们，预先讨论不同资金规则，遵守时间表，体谅截止期压力。财务计划应考虑各资助方不同的支出规定。

获批金额可能少于申请，成本也可能在获批后改变。通常，即使工作比预期贵，总额也不能增加；跨人员、设备、耗材、差旅等类别调款，或大幅减少关键人员投入比例，需资助方批准。机构的资助后管理人员负责确保基金和合同恰当使用，应定期报告、预测，协助充分利用经费，避免最终超支。少量结余可能有益，因为资助者常允许“不追加经费延期”，用余额完成或扩展批准范围的工作。
\sourcepage{568}
药企或其他私人、行业资助的预算通常属于合同，内含方案以及双方任务。合同是约束研究者和机构的法律文件，常规定达到招募目标、提交进度报告等交付内容对应的付款时间金额。需法律帮助保护知识产权、数据访问和发表权；但律师通常不熟悉具体研究工作，研究者对范围和交付内容的意见不可缺少。

\subsection{资料收集平台}
规划的重要一步是选择收集资料、实施干预的合适平台，平衡科学方法的周密性和精度、参与便利与可及性、人员成本和负担。过去需面对面评估的许多测量，如今可用网络问卷、可穿戴设备、远程医疗或视频检查。

但开发测试远程程序可能需大量时间，包括详细说明如何访问平台、使用归还借用设备，以及确保远程数据安全和质量的人员程序。知情同意还需符合电子记录监管要求，如 FDA 的21 CFR Part 11。
\sourcepage{569}
即使收集资料不必面对面，早期接触仍可能建立关系，尤其需长期投入的研究。初次访视见过人员后，参与者可能更有联系感、更愿投入，从而更可能完成无需再见面的后续活动。

\subsection{伦理审查批准}
多数研究在招募前需取得方案、同意书及招募材料批准（第7章）。面向参与者的问卷、表格、网站等也通常需先审查。应熟悉当地委员会要求和耗时；若成员提出疑问或要求修改，时间可能很长。为减少延误，应及早与委员会工作人员讨论影响参与者的程序、设计决定。

\subsection{支付覆盖分析}
影像、验血等，除主要为研究进行者外，可能由参与者保险支付。多数机构要求独立的\term{支付覆盖分析}（coverage analysis），判断服务应由保险还是资助方付费。例如，新抗癌药研究需测肿瘤大小变化，分析将判断所需影像能否向患者保险收费。

\subsection{操作手册与表格开发}
\sourcepage{570}
可扩展方案形成操作手册，纳入方案、组织政策和更具体的方法（附录18A），规定如何招募入组及每次访视活动。随机试验可包括随机化、盲法、各变量测量、质控、数据管理、分析计划和数据安全监察（第11章）。还应列出所有问卷测量，以及联系参与者、访谈、填写编码、录入编辑、采集处理标本的说明。

\textbf{多地点合作或长期研究尤其需要操作手册，}无论地点和人员如何更换，都能保持程序一致。即使一人完成，书面的操作性定义仍能减少随机变异和技术随时间改变，迫使研究者想清细节；编写时常发现此前未注意的组织困难和修改需要。

资料表设计影响质量和成功（第19章），首人入组前应全部预试。需主观判断的录入必须有明确操作定义，表上简述、手册详述。预试保证含义清楚、易用；每页标填写或访视日期、参与者及填表人员 ID，可维护完整性。网络表格、平板等移动设备也需启动期预试，并在手册中说明用法。
\sourcepage{571}
\subsection{数据库设计}
首人入组前，应开发测试用于录入、编辑、保存、监测和分析的数据库（第19章）。即使招到合适人员，开发测试仍需数周至数月。机构常提供协助，大型研究可用专业服务，但须预留预算。初学者可咨询可信的内部技术专家和资深顾问。

急于开始研究的人，有时先把数据记在纸质表格或电子表格中，而不使用真正的数据库程序。这种做法起初较容易，但到了清理和分析数据时，可能耗费更多时间与精力。及早建立数据库的优点，是促使研究者考虑每个变量的哪些值可以接受，并禁止或提示越界、不合逻辑或缺失的值。高质量的数据录入与管理系统既能改善收集和录入时的质量控制，也能减少清理所需时间。其最大价值，是避免到研究后期才发现大量缺失、越界或不合逻辑的值，而那时已经无法纠正。

\subsection{招募与留存}
许多研究最困难的是按时招募，必须及早规划时间、人员、资源、经费和专长。研究者通常低估而很少高估招募所需投入（第3章）。
\sourcepage{572}
开始后，应监测各策略的成功程度，包括每单位成本带来的合格人数，可停用部分渠道或重新分配经费。有时询问拒绝者也能指导今后策略。

前瞻性随访还需促进并监测留存。先制定减少退出的方案，频繁联系以维持兴趣，但避免负担过重。入组有奖励，随访也应有额外奖励。若留存低于预期，可增加或改变联系、调整奖励、删除繁重程序。其他策略包括：
\begin{itemize}
\item 提供停车或交通券；
\item 定期反馈骨密度或问卷汇总等个人测量结果；
\item 寄生日、节日卡或研究阶段成果简报；
\item 长而费力的访视提供食物饮料；
\item 提供晚间、周末等灵活时间。
\end{itemize}

\section{方案定稿}
\subsection{预试、完整演练与先导研究}
这些活动评价方法的可行性、效率、成本，测量可重复性和准确度，以及可能招募率。
\sourcepage{573}
规模取决于设计和需要；多数研究预试程序即可，大而昂贵的研究可能需完整先导研究。花费最终预算最多10\%作先导，以确认招募可行、样本量现实、测量适当并尽量降低负担，可能值得。

\term{预试}（pretests）由人员评价具体问卷、测量或程序的功能、适当性与可行性。例如，故意填写、录入缺失、越界或不合逻辑数据，确认数据库编辑系统能识别。

开始前应进行\term{完整演练}（dress rehearsal），让 PI 或人员完整模拟访视，解决最终工具和程序的问题。纸面上流畅的方案，实践常暴露组织和实质问题，演练可促成改进。

\term{先导研究}（pilot studies）先获取信息，指导正式研究设计实施，常决定其成功（临床试验见第11章）。它检验招募、随机分配、测量、收集、留存等可行性，并估算全部成本。

\subsection{资料收集开始后的小幅修订}
再周密的设计和预试，开始后仍常出问题。一般原则是招募开始后尽量少改。
\sourcepage{574}
但修改有时能加强研究；需权衡方法改进的收益与一致性降低、实施时间金钱及团队混乱的代价。

明确操作性定义相对容易。例如，排除酒精滥用者时，已戒酒数年者可否纳入？应与共同研究者讨论，再以备忘录沟通并记入手册，确保之后所有人员一致执行。

这类小调整常无需伦理批准，尤其未增加风险、未改变已批方案时；不确定应询问委员会。任何方案、同意书、手册或文件改变，都应标新版本号与日期，并建立防止使用过期资料的系统。

\subsection{资料收集开始后的重大修订}
改变资格、干预或结局等属于重大问题。即使理由充分，也必须考虑：若有利于解释，分别分析和报告改变前后资料。

Raloxifene Use for The Heart（RUTH）试验可说明这种判断。该多中心试验在10,101名冠心病高危女性中评价雷洛昔芬对冠脉事件的影响，最初主要结局为非致死性心梗或冠脉死亡。
\sourcepage{575}
早期发生率低于预期，可能因为溶栓和经皮血管成形等新治疗降低心梗风险。执行委员会慎重考虑后，把心梗以外的急性冠脉综合征也纳入。修订较早，之后可一致判定；此前也收集了足以按新标准判断的心脏事件资料，故可搜索既往事件（1）。

同样在 RUTH 试验早期，雷洛昔芬多结局评价（Multiple Outcomes of Raloxifene Evaluation，MORE）试验显示，雷洛昔芬治疗显著降低乳腺癌风险（2）。但乳腺癌例数少，结果尚不能定论；而且 MORE 纳入的女性全部患有骨质疏松，对结果的可推广性也有疑虑。为了确定雷洛昔芬能否同样降低无骨质疏松老年女性的乳腺癌风险，RUTH 执行委员会决定把乳腺癌列为第二个主要结局（1）。

两次改变都需修订方案、各中心伦理和 FDA 批准，并更新许多表格文件，增强研究而未损害整体完整性。重大修订只能在团队与数据安全监察委员会、资助者等适当顾问权衡后实施，分析和结论中还须处理影响。

\subsection{研究结束}
所有纵向研究终会停止随访。参与者完成最后访视的\term{结束阶段}（closeout），需仔细规划（3）。
\sourcepage{576}
至少应感谢时间和努力，告知其参与对成功至关重要，还可包括：
\begin{itemize}
\item 最后访视当面并书面，或以后邮寄，告知参与者及通常其医生有临床意义的检测结果；
\item 盲法试验在最后访视，或全部完成且主要分析完成／发表时，告知干预分组；
\item 按监管或方案要求归档纸质、电子数据；
\item 分析完成后关闭伦理审查记录；
\item 安排保存可供未来或附属研究的标本；
\item 邮寄主要论文及通俗新闻稿或结果介绍，并提供答疑电话。
\end{itemize}

\section{研究过程中的质量控制}
\subsection{药物临床试验质量管理规范}
保证各方面高质量，是临床研究的关键。\term{药物临床试验质量管理规范}（GCP）最初用于需要 FDA 等批准的药物试验，被定义为“设计、实施、记录和报告有人类受试者参与的试验的国际伦理与科学质量标准”，向公众保证参与者权利、安全与福祉受保护（4）。这些原则逐渐用于公共机构资助的各种试验及其他设计（表18.2）。
\sourcepage{577}
FDA《联邦法规》第21篇详列要求（4，5）；国际协调会议（International Conference on Harmonization）（6）提供欧美日监管机构使用的质控指南。

\begin{bookbox}{表18.2\quad GCP 涵盖的临床研究实施方面}
\begin{itemize}
\item 设计有适当的临床前、动物和其他资料支持；
\item 遵循研究伦理原则；
\item 认真执行书面方案；
\item 研究者及提供临床照护者受过训练并具备资格；
\item 临床、实验室程序达到质量标准；
\item 数据可靠、准确；
\item 保存完整、准确的记录；
\item 预先规定统计方法并认真遵循；
\item 清楚、公平地报告结果。
\end{itemize}
\end{bookbox}
最佳实施方式是为各环节制定清楚、详细的书面\term{标准操作规程}（standard operating procedures，SOP）。方案、手册、分析计划和数据安全监察计划均可视为 SOP，但常不包括培训认证、数据库开发测试、文件维护保密备份等。机构专员可提供模板。本章讨论程序及管理质控，伦理见第7章。

\subsection{临床程序的质量控制}
应指定一名团队成员担任\term{监管事务经理}或\term{质量协调员}（regulatory affairs manager or quality coordinator），负责在研究各方面实施适当的质量控制方法，包括监督人员培训与认证；
\sourcepage{578}
保存人员认证、人员 ID 和职责授权日志；准备监管申报材料；以及在研究期间监测质量控制程序的执行。目标是在可能的问题发生之前加以发现和预防。质量控制协调员还可能负责为伦理审查委员会、FDA 或资助方的审计作准备，并担任联络人。质量控制始于规划阶段，贯穿整个研究（表18.3）。

\begin{bookbox}{表18.3\quad 临床程序的质量控制\textsuperscript{a}}
\textbf{研究前：}开发操作手册；确定招募策略；建立测量操作性定义；制作标准工具和表格；建立质控体系；制定参与者和研究者盲法；指定协调员；培训并记录；认证并记录。

\textbf{研究中：}持续而关切地领导；定期团队会议；制定药物干预专门程序；重新认证；定期考评；跨技术员及随时间比较测量。

\small\textsuperscript{a}临床程序包括血压测量、结构式访谈、病历审阅等。
\end{bookbox}
\begin{itemize}
\item \textbf{操作手册。}是质控的必要条件（附录18A）。例如，身高作骨质疏松预测变量时，应规定身高计品牌、型号，准备（脱鞋袜）、站位、测量和记录。
\sourcepage{579}
\item \textbf{校准、培训与认证。}体重秤、身高计、影像和实验室设备应研究前及期间定期校准。人员应开始前受训、确认操作能力，期间按计划重新认证，并在现场保存全部记录。
\item \textbf{考评与观察。}征得参与者许可，主管应旁观每位成员各类访谈、操作的至少一个完整例子，包括代表性访视或电话。事先提供、基于方案手册的标准清单可指导观察；之后以积极而非贬低方式讨论问题、解决，并将时间结果记入培训日志。

同事互评有助士气、协作和一致执行；所有人会对质控产生主人翁意识，观察者也常与被评者学到同样多。
\item \textbf{定期数据报告。}定期汇总程序与测量技术质量，可提示缺失、错误或变异。
\sourcepage{580}
例如，过去2月血压筛查人员均值不同，可能揭示技术差异；几个月内读数标准差渐变，可能提示技术改变。报告还应涵盖招募、录入及时性、缺失越界比例、数据质疑处理时间、随访成功及干预依从性。
\item \textbf{药物专门程序。}注意标签、运输、储存、发药、回收处置。与厂家或研究药房规划并监督配送，检测盲态药物成分，确保药物剂量正确；用明确程序和日志追踪接收、储存、发放与退还。
\end{itemize}

\subsection{实验室程序的质量控制}
可采用表18.3许多方法，但标本易被错误处理或贴错标签，技术性质也要求特别策略：
\begin{itemize}
\item \textbf{标签。}把标本误标成他人 ID，之后可能无法发现或纠正。应每次采集核查姓名和 ID，使用带条形码／二维码的预印标签，避免手写误差和数字颠倒。
\sourcepage{581}
\item \textbf{盲法。}对标本而非本人测量时更易设盲；标签不应暴露分组或其他关键变量。即使自动血糖等看似客观检查，也可减少偏倚机会。但须建立异常报告渠道，由有资格者审阅、决定是否通知参与者或行动；临床试验还须安排异常提示需立即处理时的紧急揭盲。
\item \textbf{盲态重复标本、标准混合标本与共识测量。}把标本或影像送中央实验室作化学分析或判读时，可考虑通过同一系统送入盲态重复样本：从随机抽取的部分参与者另取第二份标本，赋予独立的虚构 ID。这样可估计测量技术的精密度。另一方法是在研究开始时制备一批混合标本，冷冻保存，定期取分装标本、标上虚构 ID 后送检。起初用当时可用的最佳技术测定混合标本的数值，之后在研究期间将其作为金标准，从而估计准确度和精密度。第三种方法适用于宫颈涂片或乳腺影像判读等本身存在变异的测量，由两名相互独立、设盲的读片者评价。如果两者在预定界限内一致，就确定结果；不一致时，可以讨论达成共识，也可以征求第三名读片者的意见。
\item \textbf{商业实验室合同。}有些研究使用商业实验室测定血液、血清或组织样本中的数值。实验室必须持有执照并通过认证，研究办公室应保存认证文件副本。实验室应提供变异系数等测量可重复性资料，保证及时服务，并提供处理编码标本、向研究者报告异常结果，以及把数据传入主数据库的标准程序。
\sourcepage{582}
\end{itemize}

\subsection{数据管理的质量控制}
研究开始前，应建立并预试数据管理系统（第19章）。工作包括设计记录测量结果的表格；选择用于数据录入、编辑和管理的计算机硬件与软件；为缺失、越界和不合逻辑的录入设置编辑参数；以及规划空白统计表，确认收集的变量能够满足分析需要（表18.4）。
\begin{bookbox}{表18.4\quad 数据管理质量控制：研究前的步骤}
\begin{itemize}
\item 精简，只收集所需变量；
\item 选择合适的数据库硬件软件；
\item 编程标记缺失、越界及不合逻辑值，并用这些数据测试；
\item 规划分析，用空表检验；
\item 纸质或电子表格应：含义清楚；内容一致（如选项穷尽且互斥）；录入格式清晰、箭头指示跳题；英文采用小写，必要时用大写、下划线和粗体强调；美观易读；经过预试验证（见第15章）；每页标日期、研究 ID 和／或条形码。
\end{itemize}
\end{bookbox}
\begin{itemize}
\item \textbf{缺失数据。}比例大时可造成灾难，少量有时也能产生偏倚。例如，一项操作有5\%迟发死亡，若10\%失访且死亡是常见原因，就可能严重低估并发症。
\sourcepage{583}
缺失数据造成的错误结论，有时可以事后纠正，例如在本例中全力寻找失访参与者，但测量往往无法补做。虽然统计方法可以根据基线、其他随访访视的信息，或其他参与者的均值来插补缺失值，但缺失观察很多时，不能保证结论免受未应答偏倚影响。唯一好的解决办法，是在设计和实施时尽量减少缺失，例如由团队成员在参与者完成访视前检查表格完整性，设计不允许跳过条目的电子录入界面，以及让数据库立即标出缺失值、提醒人员处理（表18.5）。应在参与者仍在诊所时处理缺失的临床测量，那时发现错误后容易纠正。
\item \textbf{不准确、不精密数据。}问题隐蔽，多人测量时尤其易漏察。持续错误技术可能造成严重偏倚，研究者却以为变量仍代表原定含义，作出错误结论。

培训认证、定期考评、比较不同人员数据均值和范围，有助识别预防。还可互动编辑，提示或拒绝缺失、不一致和越界表格。原始数据更改应有标准程序，尽早进行，并保留\term{审计轨迹}（audit trail），如划掉而不擦除原值，再签名、注明日期；电子系统也应如此。
\sourcepage{584}
它说明修改理由，也有助防止造假。定期检查重要变量频数分布，可在仍能联系参与者、请其返回时纠正旧错，并防止后续重犯。
\item \textbf{伪造数据。}负责人应警惕不诚实同事或雇员编造资料。应谨慎选人，建立密切关系促进伦理行为，审查数据时留意造假，并不定期核查原始来源，确认真实。
\end{itemize}

\begin{bookbox}{表18.5\quad 数据管理质量控制：研究中的步骤}
\begin{itemize}
\item 参与者尚在现场时标记、检查遗漏和重大错误；
\item 各页 ID 与日期无错误或颠倒；
\item 该次访视应填的全部正确表格均已完成；
\item 无漏填或错误跳题；字迹清楚；
\item 关键变量在允许范围内，彼此一致（如年龄与出生日期）；
\item 定期计算频数分布和变异指标，发现异常；
\item 制作其他定期汇总表，寻找错误。
\end{itemize}
\end{bookbox}

\subsection{多中心合作研究}
许多问题需要超过单中心可提供的参与者，因而多地合作。
\sourcepage{585}
有时地点均在一城一州，一人可监督；但常相隔数千英里，各有经费、行政与监管结构。

需专门措施保证程序一致、数据可比，才能合并分析。协调中心建立通信网络，协调手册、表格及其他标准化质控，培训各中心测量或干预人员，监督管理、分析和发表。通常采用互联网连接的分布式电子录入系统。

治理可由各 PI、资助机构代表组成指导委员会，下设小组。质控组负责标准化、培训、认证和考评，可能需集中培训、现场考评，以及协调中心和同行审计，复杂且昂贵。其他组可监督招募临床活动、审批发表报告、审议附属研究。

操作定义和方法改变常由某中心提问、相关人员或委员会回答而产生，应在网站或共享文件持续公布。积累较多时，更新手册等文件，标日期；单中心小研究可简单记录带日期说明，存入手册。
\sourcepage{586}
多中心还可能有临床研究监察员，即临床监查员，监督方案和监管合规，可直接为资助者工作，也可独立或受聘咨询。通常不收集资料，但现场访视、审查记录，与日常管理人员或 CRC 沟通。

\subsection{最后一点思考}
常见错误是收集过多数据，想纳入所有可能有趣的开始时变量，再增加随访。除时间成本，还会使参与者疲倦、厌烦甚至退出，损害更重要测量质量；数据库更大更复杂，也使质控和分析困难。

\textbf{应逐一质疑每个变量的必要性，删除许多可有可无者。}少量有意重复可改善关键变量效度，但精简应是原则。

\section{小结}
\begin{enumerate}
\item 成功实施始于为启动和研究配置人员、场地与经费，需要 PI 有力领导。
\item 启动包括预算、伦理批准，并通过预试招募、干预、预测和结局测量、表格及数据库的适当性与可行性，使方案和手册定稿，减少收集开始后的修改。
\sourcepage{587}
\item 增加问卷项、明确操作定义等小改动较易，但有时需伦理批准，也可能影响分析。
\item 干预、纳入标准或主要结局等大改动影响深远，应慎重考虑，并取得 DSMB、伦理委员会和资助方等批准。
\item 结束程序应适当告知发现，管理照护过渡及其影响。
\item 质控应由协调员监督，依 GCP 系统开展：操作手册等 SOP、培训认证考评、招募访视测量报告、定期会议；实验室标本标签与盲法、标准混合标本、盲态重复和共识测量；通过表格电子系统监督收集、录入、编辑分析的完整准确；多中心则设小组及分布式管理质控系统。
\end{enumerate}

\begin{references}
\item Mosca L, Barrett-Connor E, Wenger NK, et al. Design and methods of the Raloxifene Use for The Heart (RUTH) Study. Am J Cardiol. 2001;88:392-395.
\item MORE Investigators. The effect of raloxifene on risk of breast cancer in postmenopausal women: results from the MORE randomized trial. Multiple Outcomes of Raloxifene Evaluation. JAMA. 1999;281:2189-2197.
\sourcepage{588}
\item Shepherd R, Macer JL, Grady D. Planning for closeout—from day one. Contemp Clin Trials. 2008;29:136-139.
\item U.S. Food and Drug Administration. CFR - Code of Federal Regulations Title 21. \url{https://www.accessdata.fda.gov/scripts/cdrh/cfdocs/cfcfr/cfrsearch.cfm}
\item U.S. Food and Drug Administration. Good Clinical Practise. \url{https://www.fda.gov/about-fda/center-drug-evaluation-and-research-cder/good-clinical-practice}.
\item European Medicines Agency. Good clinical practice. \url{https://www.ema.europa.eu/en/human-regulatory/research-development/compliance/good-clinical-practice}
\end{references}

\chapterappendix{18A}{随机试验操作手册目录示例}
\sourcepage{589}
\begin{description}[leftmargin=1em,style=nextline]
\item[第1章\quad 研究方案]
\item[第2章\quad 组织与政策]2.1 参与单位（临床中心、实验室、协调中心等）、研究者和人员管理、治理（委员会、资助机构、数据安全监察等）；2.2 政策指南（论文与报告、附属研究、利益冲突等）。
\item[第3章\quad 招募]3.1 资格与排除标准；3.2 抽样设计；3.3 招募方法（宣传、转介联系、筛选等）；3.4 知情同意。
\item[第4章\quad 临床访视]4.1 首次访视内容；4.2 随访内容与时间；4.3 未应答者随访程序。
\item[第5章\quad 随机化与盲法程序]
\item[第6章\quad 预测变量]6.1 测量程序；6.2 干预，包括药物标签、配送和处理；6.3 依从性评估。
\item[第7章\quad 结局变量]7.1 主要结局评估与裁定；7.2 其他结局和不良事件评估、管理。
\item[第8章\quad 质量控制]8.1 概述与职责；8.2 程序培训；8.3 人员认证；8.4 设备维护；8.5 同行评审与现场访视；8.6 定期报告。
\item[第9章\quad 数据管理]9.1 收集与记录；9.2 录入；9.3 编辑、储存与备份；9.4 保密。
\item[第10章\quad 数据分析计划]
\sourcepage{590}
\item[第11章\quad 数据与安全监察指南]
\item[附录]A.1 致参与者、初级医疗服务者等的信件；A.2 问卷、表格；A.3 程序和标准等细节；A.4 招募材料（广告、传单、信件等）。
\end{description}
{\noindent\small 此为大型多中心试验示例，小型研究可较简略。\par}

\chapterappendix{18B}{第18章习题：研究实施与质量控制}
\sourcepage{591}
\begin{enumerate}
\item 一项研究询问“心梗住院后死亡有哪些预测因素？”助理详细提取病历，纸表记录，并广泛访谈120名住院患者，随访1年。约15\%死亡。收集完成后，一名助理用电子表格录入。开始分析才发现某些预测变量缺失10\%--20\%，另一些不合常理。仅57\%完成1年随访，有些现在已逾期一年多。请你咨询：
\begin{enumerate}\item 现在怎样改善数据质量？\item 下次至少可用哪三种方式减少缺失与错误？\end{enumerate}
\item 你正申请单中心随机试验经费，研究改善老年失眠者睡眠的新药。团队助理曾合作其他主题、干预，但没有失眠经验。申请要求时间表，说明完成主要任务所需时间。第一年草案包含下列活动；还有哪些重要活动需纳入规划？
\end{enumerate}
\begin{center}
\begin{tikzpicture}[x=2.2cm,y=.9cm,font=\small]
\fill[gray!25] (0,0) rectangle (6,.6);
\foreach \m in {1,...,6} {\node at (\m-.5,.3){第\m 月};}
\foreach \x in {0,.25,...,6} {\draw[gray!45,thin] (\x,0)--(\x,-3);}
\fill[dcrteal!35] (0,0) rectangle (2,-1);
\fill[dcrteal!35] (0,-1) rectangle (1.5,-2);
\fill[dcrteal!35] (2,-1) rectangle (3.5,-2);
\fill[dcrteal!35] (1.5,-2) rectangle (3,-3);
\node[align=center,text width=4.2cm] at (1,-.5){确定研究方案和知情同意};
\node[align=center,text width=3cm] at (.75,-1.5){确认伦理及数据安全监察批准};
\node[align=center,text width=3cm] at (2.75,-1.5){伦理修订获最终批准};
\node[align=center,text width=3cm,font=\fontsize{8.5}{10}\selectfont] at (2.25,-2.5){完成\\\mbox{ClinicalTrials.gov} 注册};
\foreach \x in {0,...,6} {\draw (\x,.6)--(\x,0);}
\foreach \y in {.6,0,-1,-2,-3} {\draw (0,\y)--(6,\y);}
\draw (0,.6)--(0,-3) (6,.6)--(6,-3);
\end{tikzpicture}
\end{center}
